\documentclass[10pt,a4paper]{amsart}
\usepackage[margin=19mm]{geometry}
\usepackage{amsmath,amssymb,mathtools}
\usepackage[scr=rsfs,cal=euler]{mathalfa}
\usepackage{xfrac}
\usepackage[unicode,hidelinks,bookmarks=false]{hyperref}
\newcommand{\ie}{i.e.\ }
\newcommand{\cf}{cf.\ }
\newcommand{\resp}{respectively\ }
\newcommand{\Subsection}{\subsection}
\providecommand{\nobreakdash}{\nobreak\hbox{-}}
\setlength{\emergencystretch}{2em}
\multlinegap0pt
\usepackage{algorithm}\usepackage{algpseudocode}
\usepackage{amssymb}
\usepackage{mathtools}
\usepackage[shortlabels]{enumitem}
\newtheorem{corollary}{Corollary}
\newtheorem{lemma}{Lemma}
\usepackage{cleveref}
\crefname{corollary}{Corollary}{Corollarys}
\crefname{lemma}{Lemma}{Lemmas}
\newcommand*{\NN}{\mathbb{N}}
\newcommand*{\RR}{\mathbb{R}}
\newcommand{\innp}[2]{\langle #1, #2 \rangle}
\title{On the Iterate Convergence of AdaGrad for Generalized Smooth Convex Optimization}
\author{Mathieu Besan\c{c}on, Tung Quoc Le}
\date{Source: arXiv:2608.00482v1 (CC BY 4.0)}
\begin{document}
\maketitle

\noindent\emph{Source and licence.} On the Iterate Convergence of AdaGrad for Generalized Smooth Convex Optimization. Version arXiv:2608.00482v1 (CC BY 4.0), distributed by arXiv under the Creative Commons Attribution 4.0 International licence, http://creativecommons.org/licenses/by/4.0/, as recorded in the arXiv licence metadata for this exact version and shown on the abstract page https://arxiv.org/abs/2608.00482v1. Full licence notice: \texttt{licenses/arxiv-2608.00482-CC-BY-4.0.txt}.\par\smallskip

\noindent\textit{Editorial scope.} Counted unit: the article's lemma lemma:dynamics-angle (source0.tex lines 1667--1674). The article's complete original proof of it is delivered with it (source0.tex lines 1675--1782, one \texttt{proof} environment of 108 lines). No mathematics is written, repaired or re-derived: the statement and the proof are the article's own lines, in the article's own notation, and no retained line is substituted or re-flowed.  Retained uncounted as the source-local context the proof needs: lines 1432--1434; lines 1476--1492; lines 1534--1565; lines 1591--1605.  Nothing was omitted from any retained span, so the package carries no omission record.  No retained line contains a citation, so the wrapper carries no bibliography and no bibliography entry.  No retained line contains a figure environment, an image-inclusion command or drawing code, so no image is delivered and the package declares no asset dependency.  Editorial formatting only, in addition: an AMS \texttt{amsart} preamble that restates the article's own declarations and macros used by the retained lines, namely the environments \texttt{corollary}, \texttt{lemma} on separate counters, together with the standard packages those lines need.  The retained environments print in this document as Corollary 1, Lemma 1 in order of appearance, whereas the article prints the same items with the numbering of its own full document.  The complete native source of the article as distributed in the arXiv e-print for this exact version is bundled unchanged as \texttt{sources/206545/source0.tex}.
\begin{corollary}[L-smoothness extension]\label{cor:L-smoothness}
    There exists an $L$-smooth function $f: \RR^2 \to \RR$ that agrees with $\tilde{f}$ on its domain, i.e., $\RR^2 \setminus B\left(0,\frac{1}{2}\right)$. 
\end{corollary}

\begin{lemma}[Choice of step-size $\eta$ and the constant $\delta$]
    \label{lemma:step-size-choice}
    If we choose $\eta, \delta > 0$:
    \begin{equation*}
        \begin{aligned}
            \frac{\eta}{\delta} \leq \frac{1}{L},    
        \end{aligned} 
    \end{equation*}
    where $L$ is the smoothness constant of the function $f$ (given in \Cref{cor:L-smoothness}), then we have:
    \begin{equation*}
        \frac{\eta}{\delta + (f(x_0) - f^\star) / \eta} I\preceq \eta \Gamma_k \preceq \frac{1}{L}I,
    \end{equation*}
    where $f^\star = \inf_{x \in \RR^2} f(x) > -\infty$. In particular, with $\eta = 1$ and $\delta \geq L$, we have:
    \begin{equation*}
        \frac{1}{\delta + f(x_0) - f^\star} I \preceq \eta \Gamma_k \preceq \frac{1}{\delta} I.
    \end{equation*}
\end{lemma}

\begin{lemma}[Dynamics of AdaGrad algorithms]
    \label{lemma:dynamics-adagrad-algos}
    Consider AdaGrad algorithms with $\eta = 1$. Let the polar coordinates of the $k$-th iteration $x_k$ be $(r_k, \theta_k)$ and define: 
    \begin{equation*}
        \begin{aligned}
            \hat{r}_k &= \begin{pmatrix}
                \cos\theta_k \\ \sin\theta_k
            \end{pmatrix}, \\
            \hat{\theta}_k &= \begin{pmatrix}
                -\sin\theta_k \\ \cos\theta_k
            \end{pmatrix},\\
            \beta_k &= \theta_k + \ln s_k.
        \end{aligned}
    \end{equation*}
    We define in addition:
    \begin{equation*}
        \begin{aligned}
            \rho_{k}^r &:= \innp{\hat{r}_k}{\Gamma_k\nabla f(x_k)},\\
            \rho_{k}^\theta &:= \innp{\hat{\theta}_k}{\Gamma_k\nabla f(x_k)}.\\
        \end{aligned}
    \end{equation*}
    If $\rho_k^r < r_k$, we obtain the following dynamics: 
    \begin{equation}
        \label{eq:dynamics-s-beta}
        \begin{aligned}
            s_{k + 1} &= 1 - \sqrt{(r_k - \rho_k^r)^2 + (\rho_k^\theta)^2},\\
            \theta_{k+1} &= \theta_k - \arctan \frac{\rho_k^\theta}{r_k - \rho^r_k},\\
            \beta_{k + 1} &= \beta_k - \arctan \frac{\rho_k^\theta}{r_k - \rho^r_k} + \ln \frac{s_{k+1}}{s_{k}},\\
            \|x_{k + 1} - x_k\| &= \sqrt{(\rho_k^r)^2 + (\rho_k^\theta)^2}.
        \end{aligned}
    \end{equation}
\end{lemma}

\begin{lemma}[Trapped in the unit ball and approaching the unit circle]
    \label{lemma:trapping-approaching}
    Consider a sequence $\{x_k\}_{k \in \NN}$ generated by AdaGrad algorithms with the hyperparameters $\eta = 1, \delta > 0$ and an initialization $x_0$ such that:
    \begin{equation*}
        \theta_0 = - \ln s_0 = - \ln (1 - r_0), 
    \end{equation*}
    where $(r_0, \theta_0)$ is the polar coordinates of $x_0$.
    Then, for a sufficiently small $s_0$ and a sufficiently large $\delta$, we have:
    \begin{equation*}
        \begin{aligned}
             \forall k \in \NN, \quad s_k &> 0,\\
            \lim_{k \to \infty} s_k &= 0.\\
        \end{aligned}
    \end{equation*}
\end{lemma}

\begin{lemma}[Dynamics of the angle]
    \label{lemma:dynamics-angle}
    Under the same assumption as in \Cref{lemma:trapping-approaching}, for sufficiently large $\delta$, there exists a constant $C$ such that for sufficiently small $s_0$, :
    \begin{equation*}
        |\beta_k| \leq C s_k^{a - b + 1}, \forall k \in \NN.
    \end{equation*}
    In particular, $\lim_{k \to \infty} \beta_k = 0$. Moreover, $\lim_{k \to \infty} |\theta_k - \theta_{k+1}| = 0$.
\end{lemma}

\begin{proof}
    To understand the dynamics of $\beta$ using $s$, we consider the gradient of $f$ reformulated using $\beta$ and $s$ (instead of $\theta$ and $r$).
    \begin{equation*}
        \begin{aligned}
            &\nabla f(x) = \nabla_r f(x) \hat{r} + \frac{1}{r} \nabla_\theta f(x) \hat{\theta}\\
            =& [-(a+1)(1 - r)^{a} - b(1-r)^{b-1}(1 - \cos(\theta + \ln(1 - r))) - (1 - r)^{b - 1}\sin(\theta + \ln(1 - r))] \hat{r} \\
            +& \frac{1}{r}(1 - r)^b \sin(\theta + \ln(1-r)) \hat{\theta}\\
            =& \left[-(a+1)s^{a} - bs^{b-1}(1 - \cos(\theta + \ln s)) - s^{b - 1}\sin(\theta + \ln s)\right] \hat{r} + \frac{1}{r}\left[s_k^b \sin(\theta + \ln s)\right] \hat{\theta}\\
            =& \left[-(a+1)s^{a} - bs^{b-1}(1 - \cos\beta) - s^{b - 1}\sin(\beta)\right] \hat{r} + \frac{1}{r}\left[s^b \sin(\beta)\right] \hat{\theta}\\
            =&\underbrace{\left[-(a+1)s^{a} - bs^{b-1}(1 - \cos\beta)\right]}_{P} \hat{r} - \left[s^b \sin(\beta)\right]\underbrace{\left(\frac{1}{s}\hat{r} - \frac{1}{r}\hat{\theta}\right)}_{=:\zeta}.
        \end{aligned}
    \end{equation*}
    In the following, we will use $P_k$ and $\zeta_k$ to indicate those quantities at the $k$th iteration. 
    
    Next, we will linearize the dynamics of $\beta_k$. By \Cref{lemma:dynamics-adagrad-algos}, we have:
    \begin{equation*}
        \beta_{k + 1} = \beta_k - \arctan \frac{\rho_k^\theta}{r_k - \rho^r_k} + \ln \frac{s_{k+1}}{s_{k}},
    \end{equation*}
    where $\rho_k^r$ and $\rho_k^\theta$ are defined as in \Cref{lemma:dynamics-adagrad-algos}. We consider each term separately: If $s_0$ is sufficiently small, we have:
    \begin{equation*}
        \begin{aligned}
            \left|\arctan \frac{\rho_k^\theta}{r_k - \rho^r_k} - \frac{\rho_k^\theta}{r_k}\right| & \leq \left|\arctan \frac{\rho_k^\theta}{r_k - \rho^r_k} - \frac{\rho_k^\theta}{r_k - \rho^r_k}\right| + \left|\frac{\rho_k^\theta}{r_k - \rho^r_k} - \frac{\rho_k^\theta}{r_k}\right|\\
            &\leq \underbrace{\frac{1}{3}\left|\frac{\rho_k^\theta}{r_k - \rho^r_k}\right|^3}_{O(\|\nabla f(x_k)\|^3)} + \underbrace{\left|\frac{\rho_k^\theta\rho_k^r}{r_k(r_k - \rho^r_k)}\right|}_{O(\|\nabla f(x_k)\|^2)}.
        \end{aligned}
    \end{equation*}
    where we use the fact that $|\arctan x - x| \leq \frac{1}{3}|x|^3, \forall |x| < 1$ (using the Taylor expansion of the function $\arctan$ around $0$) and $|\rho_k^r|, |\rho_k^\theta| \leq \|\Gamma_k \nabla f(x_k)\| = O(\|\nabla f(x_k)\|)$.

    We deal with the second term (while always assuming that $s_0$ is sufficiently small) as:
    \begin{equation*}
        \begin{aligned}
            \left|\ln \frac{s_{k + 1}}{s_k} - \frac{\rho_k^r}{s_k} \right|&\leq \left|\ln \frac{s_{k + 1}}{s_k} - \frac{s_{k + 1} - s_k}{s_k} \right| + \left|\frac{s_{k + 1} - s_k - \rho^r_k}{s_k}\right|\\
            &= \left|\ln \frac{s_{k + 1}}{s_k} - \frac{s_{k + 1} - s_k}{s_k} \right| + \left|\frac{r_{k + 1} - r_k + \rho^r_k}{s_k}\right|\\
            &\leq \left(\frac{s_{k + 1} - s_k}{s_k}\right)^2 + \left|\frac{\sqrt{(r_k - \rho_k^r)^2 + (\rho_k^\theta)^2} - (r_k - \rho^r_k)}{s_k}\right|\\
            &\leq \underbrace{\left(\frac{s_{k + 1} - s_k}{s_k}\right)^2}_{=O\left(\frac{\|\nabla f(x)\|^2}{s_k^2}\right)} + \underbrace{\left|\frac{(\rho_k^\theta)^2}{s_k(r_k - \rho^r_k)\sqrt{(r_k - \rho_k^r)^2 + (\rho_k^\theta)^2}}\right|}_{=O\left(\frac{\|\nabla f(x)\|^2}{s_k}\right)},\\
        \end{aligned}
    \end{equation*}
    where we use again the fact that $|\ln(1 + x) - x| \leq x^2$ for $|x| < 1/2$ (using the Taylor expansion again). We conclude that:
    \begin{equation*}
        \beta_{k + 1} - \beta_k = -\frac{\rho_k^\theta}{r_k} + \frac{\rho_k^r}{s_k} + O\left(\frac{\|\nabla f(x_k)\|^2}{s_k^2}\right).
    \end{equation*}
    Due to the definitions of $\rho_k^\theta$ and $\rho_k^r$, we have:
    \begin{equation*}
        -\frac{\rho_k^\theta}{r_k} + \frac{\rho_k^r}{s_k} = -\frac{\innp{\hat{\theta}}{\Gamma_k\nabla f(x_k)}}{r_k} + \frac{\innp{\hat{r}}{\Gamma_k\nabla f(x_k)}}{s_k} = \innp{\zeta_k}{\Gamma_k\nabla f(x_k)}.
    \end{equation*}
    Thus, we end up with the following estimation:
    \begin{equation*}
        \begin{aligned}
            \beta_{k + 1} &= \beta_k + P_k\innp{\zeta_k}{\Gamma_k\hat{r}_k} - \sin \beta_k \underbrace{s^b\innp{\zeta_k}{\Gamma_k\zeta_k}}_{\Lambda_k} + O\left(\frac{\|\nabla f(x_k)\|^2}{s_k^2}\right)\\
            &=\beta_k (1 - \Lambda_k) + P_k\innp{\zeta_k}{\Gamma_k\hat{r}_k} - (\sin \beta_k - \beta_k)\Lambda_k + O\left(\frac{\|\nabla f(x_k)\|^2}{s_k^2}\right).
        \end{aligned}
    \end{equation*}
    We estimate the terms in the above equation: since $\lambda_- I \preceq \Gamma_k \preceq \lambda_+ I$ (with $\lambda_\pm$ determined uniquely by $\delta$ and $f(x_0) - f^\star$, see \Cref{lemma:step-size-choice}), $\zeta_k = \Theta(1/s_k)$, we have:
    \begin{equation*}
        \Lambda_k = s_k^b \innp{\zeta_k}{\Gamma_k\zeta_k} = \Theta(s_k^{b - 2}).
    \end{equation*}
    Moreover,
    \begin{equation*}
        \begin{aligned}
            |P_k\innp{\zeta_k}{\Gamma_k\hat{r}_k}| &= (a+1)|s_k^{a} \innp{\zeta_k}{\Gamma_k\hat{r}_k}| + b|s_k^{b-1}(1 - \cos\beta_k) \innp{\zeta_k}{\Gamma_k\hat{r}}| \\
            &\leq \underbrace{(a+1)|s_k^{a} \innp{\zeta_k}{\Gamma_k\hat{r}_k}|}_{= O(s_k^{a - 1})} + \underbrace{\frac{b}{2}|s_k^{b-1}\beta_k^2 \innp{\zeta_k}{\Gamma_k\hat{r}_k}|}_{=O(s_k^{b - 1}\beta_k^2)}
        \end{aligned} 
    \end{equation*}
    \begin{equation*}
        \begin{aligned}
            |(\sin(\beta_k) - \beta_k)\Lambda_k| \leq \frac{1}{6}|\beta_k^3\Lambda_k| = O(\beta_k^3s_k^{b - 2}).
        \end{aligned}
    \end{equation*}
    We also have to estimate the difference between $s_{k + 1}$ and $s_k$, which is:
    \begin{equation*}
        \begin{aligned}
            \|\underbrace{s_{k + 1} - s_k}_{\Delta_k}\| &\leq \|\Gamma_k\nabla f(x_k)\| = O(\|P_k \hat{r}_k - s_k^b \sin (\beta_k) \zeta_k\|) = O(s_k^a + s_k^{b - 2}\beta_k^2 + s_k^{b - 1}\beta_k).
        \end{aligned}
    \end{equation*}
    This estimation allows us to have:
    \begin{equation*}
        \begin{aligned}
            s_{k+1}^{a - b + 1} - s_k^{a - b + 1} &= O(s_k^{a - b}\|\Delta_k\|).
        \end{aligned}
    \end{equation*}
    Thus, we come to the following estimation: We prove by induction that $|\beta_k| \leq Cs_k^{a - b + 1}$ for a big enough constant $C > 0$ (that we will explain how to choose later). For $k = 0$, the claim holds trivially since $\beta_0 = 0$. Assume that the claim holds up to $k \in \NN$, we have the following order for each quantity that we just estimated under the assumption that $a > b - 1$:
    \begin{equation*}
        \begin{aligned}
            \Lambda_k &= \Theta(s_k^{b - 2}),\\
            |P_k\innp{\zeta_k}{\Gamma_k\hat{r}_k}| &= O(s_k^{a - 1} + s_k^{b-1 + 2(a-b+1)}) = O(s_k^{a - 1}),\\
            |(\sin(\beta_k) - \beta_k)\Gamma_k| &= O(\beta_k^3s_k^{b - 2}) = O(s_k^{3(a - b + 1) + b - 2}) = O(s_k^{3a -2b + 1}),\\
            \|s_{k+1}-s_k\| &= O(s_k^{a} + s_k^{b - 1 + 2(a - b + 1)} + s_k^{b - 1 + a - b + 1}) = O(s_k^{a}),\\
            s_{k + 1}^{a - b + 1} - s_k^{a - b + 1} &= O(s_k^{2a - b}).\\
            \frac{\|\nabla f(x_k)\|^2}{s_k^2} &= O(s_k^{2a - 2}).
        \end{aligned}
    \end{equation*}
    We choose $C > 0$ such that $C\Lambda_ks_k^{a - b + 1} - |P_k \innp{\zeta_k}{\Gamma_k\hat{r}_k}| > C's_k^{a - 1}$, for some $C' > 0$, which is always possible because $|P_k\innp{\zeta_k}{\Gamma_k\hat{r}_k}| = O(s_{k}^{a - 1})$ and $C\Gamma_ks_k^{a - b + 1} = \Theta(s_{k}^{a - 1})$. Assuming that $s_0$ is small enough so that $\Gamma_k < 1, \forall k \in \NN$, we have:

    \begin{equation*}
        \begin{aligned}
            |\beta_{k+1}| &\leq (1 - \Lambda_k)|\beta_k| + |P_k\innp{\zeta_k}{\Gamma_k\hat{r}_k}| + |(\sin \beta_k - \beta_k)\Lambda_k| + O\left(\frac{\|\nabla f(x_k)\|^2}{s_k^2}\right)\\
            &\leq (1 - \Lambda_k) Cs_k^{a - b + 1} + |P_k\innp{\zeta_k}{\Gamma_k\hat{r}_k}| + O(s_k^{3a - 2b + 1} + s_k^{2a - 2})\\
            &\leq Cs_k^{a - b + 1} - C's_k^{a - 1} + O(s_k^{3a - 2b + 1} + s_k^{2a - 2})\\
            &\leq Cs_{k+1}^{a - b + 1} - C's_k^{a - 1} + O(s_k^{3a - 2b + 1} + s_k^{2a - 2} + s_k^{2a - b}).
        \end{aligned}
    \end{equation*}
    Note that by our choice of $a > b - 1$, we have $a - 1 < \min(3a - 2b + 1, 2a - 2, 2a - b)$. Therefore, with sufficiently small $s_0$, the induction part holds.

    Finally, one has:
    \begin{equation*}
        \lim_{k \to \infty} |\theta_k - \theta_{k + 1}| = \lim_{k \to \infty}\left|\arctan \frac{\rho_k^\theta}{r_k - \rho^r_k}\right| \leq \lim_{k \to \infty} \frac{\|\Gamma_k\nabla f(x_k)\|}{|r_k - \rho^r_k|} = 0.
    \end{equation*}
    That concludes the proof.
\end{proof}

\end{document}
